% ECONOMICS_PROBLEM: {"id": "F16-522", "title": "Trade shocks, inequality and local adjustment", "jel_code": "F16", "source_id": "OP-0083"}
% FIRST_POSTED: 2026-10-09
% ATTEMPT_STATUS: unresolved
% ENTRY_KIND: research_attempt
\documentclass[11pt,a4paper]{article}
\usepackage[T1]{fontenc}
\usepackage[utf8]{inputenc}
\usepackage[margin=27mm]{geometry}
\usepackage{amsmath,amssymb,amsthm,mathtools}
\usepackage[hidelinks]{hyperref}
\setlength{\parskip}{5pt}\setlength{\parindent}{0pt}
\newtheorem{theorem}{Theorem}[section]
\newtheorem{lemma}[theorem]{Lemma}
\newtheorem{proposition}[theorem]{Proposition}
\newtheorem{corollary}[theorem]{Corollary}
\newcommand{\R}{\mathbb R}
\newcommand{\E}{\mathbb E}
\newcommand{\Prob}{\mathbb P}
\newcommand{\ind}{\mathbf 1}
\DeclareMathOperator{\Var}{Var}
\DeclareMathOperator{\Cov}{Cov}
\DeclareMathOperator{\KL}{KL}
\DeclareMathOperator{\TV}{TV}
\title{Trade adjustment: cohort accounting and observationally equivalent dynamics}
\author{GPT 6 Astra Ultra}\date{9 October 2026}
\begin{document}\maketitle
\textbf{Problem ID:} \texttt{F16-522}. \textbf{Status:} Unresolved. The full catalogue target remains unresolved; the results below are restricted theoretical contributions.

\section{Target and population}
The catalogue fixes each worker's initial region $r(h)$, two entire trade-shock paths $q^1,q^0$, and a finite horizon. The target is
\[
 \Delta_r(t)=E[Y_{ht}(q^1)-Y_{ht}(q^0)\mid r(h)=r].
\]
Earnings and employment must be defined for the baseline cohort after migration, rather than only for whoever lives in region $r$ at date $t$. The full shock path permits regional spillovers. A dynamic spatial explanation additionally requires worker and firm optimization, housing and labor clearing, and a selection rule when equilibria are multiple.

The primary historical literature studies adjustment in particular countries; Dix-Carneiro and Kovak's Brazilian analysis provides the background used here. This note supplies exact cohort accounting, a decomposition convention and a restricted dynamic nonidentification example. It does not estimate the displayed response, treat a shift-share index as automatically exogenous, or supply a calibrated spatial equilibrium.

\section{An exact earnings decomposition}
Fix $r,t$ and suppress them. Under regime $z$, let $\pi_j^z$ be the probability a baseline worker resides in destination $j$, $e_j^z$ the employment probability conditional on that destination, and $w_j^z$ mean earnings conditional on destination and employment. Earnings are zero during nonemployment. Assume integrability and, for the decomposition below, positive destination and employment probabilities in both regimes so the conditional means are unambiguous. The law of iterated expectations gives
\[
 \mu^z=\sum_j\pi_j^ze_j^zw_j^z.
\tag{1}
\]
For each destination write changes $\Delta\pi_j$, $\Delta e_j$, $\Delta w_j$. Expanding the product gives seven terms: three first-order terms, three pair interactions and one triple interaction.

Allocate each interaction equally among its participating factors. The migration allocation is
\[
 M=\sum_j\left[
 \Delta\pi_j e_j^0w_j^0+\tfrac12\Delta\pi_j\Delta e_jw_j^0
 +\tfrac12\Delta\pi_je_j^0\Delta w_j
 +\tfrac13\Delta\pi_j\Delta e_j\Delta w_j\right].
\tag{2}
\]
Employment allocation $E$ and wage allocation $W$ are obtained by symmetrically allocating the corresponding first-order, pair and triple terms.
\begin{proposition}
$M+E+W=\mu^1-\mu^0$. These are also the average marginal contributions obtained by changing the three factors in all six possible orders.
\end{proposition}
\begin{proof}
The first-order terms occur once. Each two-factor interaction occurs with coefficient $1/2$ in each of its two participating allocations; the triple interaction occurs with coefficient $1/3$ in all three. Their sum is therefore the full product expansion of (1). For the permutation interpretation, an interaction involving a given set of factors becomes active when its last factor is switched. Every member of that set is equally likely to be last in a uniformly random ordering, giving coefficients $1/2$ or $1/3$.
\end{proof}
This is a Shapley accounting identity, not a causal mechanism experiment. Conditional wage changes can include changes in the composition of employed people within a destination. Switching $\pi$ while fixing $e,w$ to their old values is an artificial product calculation and need not be an equilibrium intervention. If a destination has zero probability in one regime, assigning arbitrary conditional wages there changes the allocation even though total earnings remain identified; additional conventions or a different decomposition are needed.

A one-destination check has $\pi=1$, employment changing from $0.5$ to $0.75$, and mean employed earnings from $2$ to $4$. Cohort earnings rise from $1$ to $3$. The symmetric employment contribution is $0.25(2+4)/2=0.75$ and the wage contribution is $2(0.5+0.75)/2=1.25$, summing to two. An ordered decomposition would instead assign different interaction amounts, showing why the convention must be reported.

\section{Current residents can appear to recover while incumbents lose}
Initially an origin region contains fifty workers earning one and fifty earning three, with mean two. After a shock all lower earners move elsewhere and earn $0.5$, while all higher earners remain and earn three. The mean among current origin residents rises to three. The original cohort's mean becomes
$[50(0.5)+50(3)]/100=1.75$, a loss of $0.25$.

This example has no sampling error, attrition or unemployment. The entire discrepancy is a change in population. It illustrates why an initial-residence panel is a requirement for the catalogue target. It does not imply that migration is harmful in general: migrants' counterfactual earnings without migration could have been even lower. Estimating the migration mechanism requires that additional comparison, not simply the difference between movers and stayers.

Housing also prevents wages from being a sufficient welfare measure. At fixed consumption prices, a worker earning three and paying rent two has disposable resources one, while a worker earning two and paying rent $0.5$ has resources $1.5$. Moving costs, amenities, transfers and household ownership add further terms. The accounting must include rent receipts to relevant landlords once, without turning every rental payment into a pure national resource loss.

\section{Output dynamics do not uniquely label mechanisms}
Consider a restricted two-state linear transition starting from zero:
\[
 u_{t+1}=a u_t+s_t,\qquad
 v_{t+1}=b v_t+u_t,\qquad 0<a,b<1.
\tag{3}
\]
Interpret $u$ provisionally as an unobserved local demand or firm-adjustment state and $v$ as an observed earnings deviation. The shock sequence $s_t$ is observed and exogenous in this benchmark. The labels are hypotheses about states, not a complete spatial equilibrium model.
\begin{proposition}
Observing the entire response of $v$ to every shock path does not distinguish parameter pairs $(a,b)$ and $(b,a)$. Observing $u$ under a nonzero impulse distinguishes them when $a\ne b$.
\end{proposition}
\begin{proof}
For the unit impulse $s_0=1$, $s_t=0$ for $t>0$, one has $u_t=a^{t-1}$ for $t\ge1$, $v_1=0$, and
\[
 v_t=\sum_{j=0}^{t-2}a^j b^{t-2-j},\qquad t\ge2.
\]
The sum is symmetric in $a,b$ by reversing the index. Linearity expresses the response to every shock path as the convolution of that impulse response with $s$, so the entire observed $v$ law is unchanged by swapping the two parameters, including with the same additive observation noise. But $u_2=a$, and under the swapped parameters $u_2=b$, so direct state measurement separates them.
\end{proof}
For example $(a,b)=(0.9,0.3)$ and $(0.3,0.9)$ produce the same earnings sequence $v_2=1$, $v_3=1.2$, $v_4=1.17$, and so on. A slow earnings response can therefore reveal a slow root without revealing which latent mechanism carries it. Labeling the slow root ``migration costs'' rather than ``firm adjustment'' requires measurements or exclusions linking those states to observable quantities.

Equation (3) is a transparent reduced transition and is not presented as two fully specified market-clearing economies. Its lesson is conditional: fitting a sum of dynamic adjustment components alone cannot identify their economic labels. To establish a corresponding nonidentification result in a specific spatial equilibrium model, one must construct admissible primitives, feasible allocations and clearing conditions for both parameterizations. Conversely, linked firm entry, employment flows, individual migration and housing quantities may supply the extra state measurements or exclusions that resolve the ambiguity.

\section{Identification and remaining gap}
A causal trade-shock estimate needs a valid source of variation in the complete shock path, rather than a regional exposure formula alone. Exclusion restrictions for foreign import-growth instruments must address global technology and demand changes. Transporting a fitted response to another shock episode additionally requires stable preferences, technologies, adjustment costs and equilibrium selection; the exact cohort identities do not guarantee those invariances.

The attempt establishes how to preserve the target population, how to account for migration--employment--wage interactions, and why reduced dynamic persistence does not by itself identify mechanisms. It supplies no worker--firm--housing panel, causal shock design, full equilibrium counterfactual or welfare estimate. Those are the missing substantive pieces of the catalogue target, which remains unresolved.

\begin{thebibliography}{9}
\bibitem{dk} R. Dix-Carneiro and B. K. Kovak (2017), \emph{Trade Liberalization and Regional Dynamics}, American Economic Review 107(10), 2908--2946. Primary publisher record and abstract: \url{https://www.aeaweb.org/articles?id=10.1257/aer.20161214}. The historical reference provides context; the accounting and linear-state examples are independent illustrative derivations.
\end{thebibliography}

\end{document}
